\documentclass[11pt]{article}
\usepackage{preamble}
\renewcommand{\answer}[1]{}

\begin{document}

\ladderheading{AP Statistics \textperiodcentered\ Unit 3 \textperiodcentered\ Topic 3.9 \textperiodcentered\ B: 2026-12-10 \textperiodcentered\ E: 2027-01-20}{When Two Proportions Are Not Independent}

\focusbanner{TODAY'S PROBLEM}

A library studies self-checkout before and after a redesign. Records give $p_B=0.42$ among 4{,}000 before visits and $p_A=0.54$ among 5{,}000 after visits.\par\smallskip \textbf{Plan I:} Draw an SRS of 120 before visits and a separate SRS of 150 after visits; do not track cardholders.\par \textbf{Plan P:} Measure the same 120 cardholders before and after.\par\smallskip \textbf{Sora:} Separate selection warrants Plan I's independent model; 0.0350 means a difference at most 0.01 is unusual.\par \textbf{Mateo:} Pairing controls person differences, so Plan P may help. Because both proportions use 120 people, Plan I's model applies to P too.\par
\begin{center}
\textbf{Sampling plans}\par\smallskip
\begin{tabular}{|l|c|c|c|}
\hline
\textbf{Plan} & \textbf{Before} & \textbf{After} & \textbf{Repeated} \\ \hline
\textbf{I} & SRS 120 & SRS 150 & no \\ \hline
\textbf{P} & 120 & same 120 & yes \\ \hline
\end{tabular}
\end{center}
\begin{firsttakebox}{FIRST TAKE --- before the video (5 min)}
Which parts of Sora's and Mateo's reasoning are warranted? Cite one design detail.
\workspace{0.9in}
\end{firsttakebox}

\begin{callout}{\IconBook\ THE VARIANCE FORMULA REQUIRES INDEPENDENT SAMPLES (keep on the page)}{calloutgreen}
For independent samples, $\mu_{\widehat p_1-\widehat p_2}=p_1-p_2$ and
$\sigma_{\widehat p_1-\widehat p_2}=\sqrt{p_1(1-p_1)/n_1+p_2(1-p_2)/n_2}$ when each sample is
below 10\% of its population. Approximate normality requires all four expected counts to be at least 10.
Repeated measurements on the same people are paired, not independent; their association belongs in the variability.
\end{callout}

\newpage
\textbf{Finish after the video:}\par\smallskip

\dokbadge{1}~\textbf{(a)} For after minus before, identify the parameter and statistic and state Plan I's sampling-distribution mean.\par
\workspace{0.8in}

\dokbadge{2}~\textbf{(b)} For Plan I, check both 10 percent conditions and all four counts. Find the standard deviation and approximate $P(\widehat p_A-\widehat p_B\leq0.01)$; interpret it.\par
\workspace{1.4in}

\dokbadge{3}~\textbf{(c)} Adjudicate the claims. Use design and conditions to place the independent model and assess Plan P. Explain why equal sample sizes are not independence, how transferring 0.0607 and 0.0350 could mislead, and what paired analysis must retain.\par
\workspace{2.0in}

\begin{sentenceframebox}\raggedright\textbf{Frame:} Plan \blank[0.4in] permits the independent model because \blank[2.5in], and its four counts are \blank[2.0in].\end{sentenceframebox}

\begin{sentenceframebox}\raggedright\textbf{Frame:} For Plan P, the missing \blank[1.2in] could make 0.0350 \blank[2.1in]; a valid analysis must keep \blank[2.0in].\end{sentenceframebox}

\begin{summaryexitbox}{EXIT --- turn this in}
One thing the video changed about my first take: \hrulefill\par
\workspace{0.4in}
\end{summaryexitbox}

\end{document}
